% \paragraph{Contexto e motivação.}
As redes de comunicação veicular passaram por algumas transições tecnológicas importantes. Inicialmente, o foco estava nas Redes Veiculares Ad hoc (\acrotip{VANET}s), que priorizavam a conectividade direta entre veículos, mas enfrentavam desafios relacionados à latência e à conectividade devido à alta mobilidade dos nós~\cite{kandali2025}. Com o amadurecimento da Internet das Coisas (\acrotip{IoT}) e a expansão de sensores embarcados na infraestrutura urbana, desenvolveu-se um paradigma de rede de grande escala que integra veículos, infraestrutura, pedestres e a nuvem, conhecido como ecossistema da Internet dos Veículos (\acrotip{IoV})~\cite{kandali2025}. Tal ecossistema foi idealizado para viabilizar os Sistemas de Transporte Inteligentes (\acrotip{ITS}) verdadeiramente autônomos e conectados, possibilitando às aplicações veiculares atender simultaneamente a requisitos de eficiência, confiabilidade e segurança~\cite{uddin2025}.

No entanto, a \acrotip{IoV} depende da capacidade de atendimento aos requisitos de baixa latência e alta vazão das tecnologias de rádio. Dessa forma, a comunicação veicular passou por algumas evoluções, partindo inicialmente da alocação de \SI{75}{\mega\hertz} na banda de \SI{5{,}9}{\giga\hertz} para o padrão Dedicated Short-Range Communications (\acrotip{DSRC}), que definiu os parâmetros para comunicação Vehicle-to-Vehicle (\acrotip{V2V}) e Vehicle-to-Infrastructure (\acrotip{V2I}) nas primeiras \acrotip{VANET}s~\cite{zheng2015}. A partir desse legado, desenvolveram-se dois eixos tecnológicos de comunicação. O primeiro, o \acrotip{IEEE}~802.11bd (Next-Generation \acrotip{V2X}), estendeu o \acrotip{IEEE}~802.11p incorporando robustez contra desvanecimento rápido e suporte a velocidades de até \SI{250}{\kilo\meter\per\hour}~\cite{naik2019}. Já o segundo direciona-se à superação das limitações do \acrotip{DSRC} em escalabilidade e custo de implantação, o que impulsionou o desenvolvimento do Cellular \acrotip{V2X} (\acrotip{CV2X}) pelo \acrotip{3GPP}~\cite{chen2017}. Sob a perspectiva regulatória, a disputa para adesão dessas tecnologias foi resolvida pela Federal Communications Commission (\acrotip{FCC}), que estabeleceu o \acrotip{CV2X} como padrão obrigatório para a comunicação de \acrotip{ITS}, com a descontinuação oficial do \acrotip{DSRC} marcada para dezembro de 2026~\cite{fcc2024_59ghz}.

A evolução das tecnologias de comunicação conseguiu atender aos requisitos de latência de rede e vazão dos \acrotip{ITS}. No entanto, as limitações em relação à capacidade computacional dos dispositivos embarcados ainda precisavam ser superadas. Nesse contexto, a Computação de Borda de Múltiplos Acessos (\acrotip{MEC}) e sua especialização veicular, a Computação de Borda Veicular (\acrotip{VEC}), viabilizam o funcionamento de aplicações sensíveis à latência ao posicionarem recursos de processamento e armazenamento próximos à rede de acesso por rádio~\cite{mao2017,miriyala2026dl_offload_review}. Dessa forma, os usuários finais podem transferir tarefas computacionalmente intensivas para os servidor de borda utilizando a infraestrutura \acrotip{5G} \acrotip{CV2X}. Essa delegação reduz drasticamente o tempo de resposta e o consumo de energia a bordo~\cite{liu2019,fang2021mec_iot,long2025}, garantindo, assim, a Qualidade de Serviço (\acrotip{QoS}) de aplicações sensíveis à latência.

Em ambientes veiculares, as capacidades computacionais das Unidades de Bordo Interconectadas (\acrotip{OBUI}s) são inerentemente heterogêneas~\cite{luo2024}. Essa diversidade de hardware, se bem aproveitada, possibilita a formação de ambientes de execução cooperativos, nos quais veículos equipados com unidades de processamento de alto desempenho disponibilizam sua capacidade ociosa temporária para auxiliar nós vizinhos com restrições computacionais e energéticas. O compartilhamento direto desses recursos melhora consideravelmente a \acrotip{QoS} global dos \acrotip{ITS}~\cite{he2019d2d_mec_twc}. Dessa forma, a integração da alta capacidade do servidor de borda com a disponibilização de processamento ocioso das \acrotip{OBUI}s fecha o ciclo de melhorias necessário para viabilizar a execução de aplicações veiculares computacionalmente intensivas e sensíveis à latência e ao consumo energético.

A operacionalização desse potencial colaborativo tem como mecanismo central o descarregamento de tarefas. Essa estratégia consiste em transferir a responsabilidade de computar tarefas complexas de dispositivos com recursos limitados para nós de processamento remoto mais poderosos, como o servidor de borda ou veículos vizinhos com maior capacidade de processamento~\cite{wang2022}. O objetivo é contornar as restrições locais de bateria, armazenamento e processamento com o intuito de reduzir o consumo de energia, minimizar a latência e melhorar a \acrotip{QoS}. O grande desafio, contudo, reside em decidir dinamicamente quando e para onde descarregar essas tarefas. Explorar plenamente esse potencial exige o desenvolvimento de políticas de alocação capazes de se adaptar à topologia altamente dinâmica da rede e de avaliar múltiplos critérios de decisão de forma simultânea e em tempo real.

O descarregamento de tarefas em redes veiculares configura-se fundamentalmente como um problema de otimização multiobjetivo~\cite{wang2020survey,chen2025}. Nesse contexto, diversas políticas têm sido propostas para solucionar esse problema. As heurísticas baseadas em regras lógicas e os métodos de Tomada de Decisão Multicritério (\acrotip{MCDM}), a exemplo das estratégias gulosas, são amplamente utilizados devido à sua simplicidade computacional e facilidade de implementação~\cite{ahmed2022survey}. No entanto, essas abordagens avaliam apenas a situação instantânea da rede, mostrando-se incapazes de antecipar as dinâmicas de tráfego e os estados futuros do ambiente. Por outro lado, as meta-heurísticas, tais como algoritmos genéticos, otimização por enxame de partículas e colônia de formigas, oferecem garantias mais sólidas de busca global e têm sido aplicadas com sucesso na tomada de decisão de descarregamento em cenários estáticos~\cite{li2020genetic,you2021efficient,kishor2022task}. Contudo, esses métodos podem se mostrar inadequados para os ambientes veiculares, visto que a busca baseada em população exige dezenas a centenas de avaliações de aptidão iterativas a cada instante de decisão, podendo causar atrasos inaceitáveis~\cite{wu2024,cao2024survey}. Essa característica de convergência relativamente lenta pode ser incompatível com os rigorosos prazos de execução na escala de milissegundos característicos das aplicações veiculares~\cite{miriyala2026dl_offload_review,uddin2025}. Além disso, esses algoritmos precisam ser reiniciados a cada alteração topológica do ambiente, e a sua complexidade de inicialização e convergência cresce significativamente com o aumento do tamanho da rede, podendo degradar a qualidade da solução em cenários de larga escala~\cite{wang2020survey,long2025}.

Cenários de avaliação simplificados, com número fixo de veículos, distribuições de tarefas uniformes e modelos de canal idealizados, não capturam a complexidade e o alto dinamismo das redes veiculares reais~\cite{miriyala2026dl_offload_review}. Políticas de descarregamento treinadas e avaliadas nesses ambientes podem sofrer degradação de desempenho quando a densidade veicular ou a distribuição da carga de trabalho muda repentinamente~\cite{miriyala2026dl_offload_review}, revelando uma dicotomia crítica entre o sucesso obtido em simulações e a real aplicabilidade em cenários operacionais. Nesse contexto, é fundamental que o desenvolvimento de políticas de descarregamento robustas busque ativamente mitigar essa disparidade. Isso exige a concepção de representações de estado adaptativas que garantam a capacidade de generalização do modelo perante as constantes alterações no número de nós candidatos ao descarregamento  eliminando a dependência de ambientes de tráfego estacionários~\cite{shao2025,uddin2025}.

O \acrotip{DRL} apresenta-se como uma solução natural para o descarregamento de tarefas porque aprende uma política de controle que mapeia observações em ações considerando as consequências futuras de longo prazo, e não apenas o estado instantâneo da rede~\cite{liu2019,shi2023}. Diferentemente das heurísticas e das meta-heurísticas, um agente \acrotip{DRL} possui a capacidade de explorar correlações temporais nos padrões de tráfego, na dinâmica das filas, na mobilidade veicular e na qualidade do canal de comunicação a partir da interação contínua com o ambiente~\cite{li2025drl_survey,uddin2025}. 

No entanto, as arquiteturas \acrotip{DRL} padrão compartilham uma fraqueza estrutural crítica relacionada à sua aplicabilidade prática, pois frequentemente exigem retreinamento sempre que o cenário operacional sofre mudanças na quantidade de nós computacionais candidatos~\cite{miriyala2026dl_offload_review}. Como a camada de entrada das redes neurais profundas depende de um vetor de tamanho fixo, as constantes alterações no número de nós candidatos ao descarregamento, como o servidor de borda e veículos vizinhos, acabam invalidando o modelo previamente treinado~\cite{uddin2025}.

% Em trabalho anterior, propusemos o FOFF~\cite{vieira2025foff}, um framework Fuzzy-TOPSIS que ranqueia o servidor de borda e veículos vizinhos em cenários \acrotip{VEC} estáticos usando pesos linguísticos de critérios para tratar incertezas. O FOFF demonstrou que métodos multicritério podem alcançar forte eficiência energética e \acrotip{QoS} sem qualquer fase de treinamento. Sua limitação é precisamente a falta de adaptação, visto que o Fuzzy-TOPSIS seleciona o melhor candidato por uma regra de pontuação fixa e não pode melhorar com a experiência ou ajustar-se a tráfego não estacionário. Com base nesses aprendizados, propomos o \acrotip{SAC}-\acrotip{TOPSIS}, uma arquitetura híbrida que delega a avaliação multicritério de candidatos ao \acrotip{TOPSIS} e a otimização de política de longo prazo ao algoritmo \acrotip{SAC}, combinando a otimalidade local do \acrotip{MCDM} com a capacidade prospectiva do \acrotip{DRL}.

Neste contexto, esse trabalho propõe uma política híbrida de descarregamento de tarefas denominada \acrotip{SAC}-\acrotip{TOPSIS} capaz de lidar com as restrições de funcionamento das \acrotip{DRL} em um ambiente veicular dinâmico. Em resumo, as principais contribuições são:

\begin{enumerate}[leftmargin=*,nosep]
  \item Desenvolvimento de uma solução inovadora para o descarregamento de tarefas que combina a técnica \acrotip{TOPSIS} com o algoritmo \acrotip{SAC}. O \acrotip{TOPSIS} fornece a melhor alternativa baseada nas condições instantâneas da rede, enquanto o \acrotip{SAC} aprende a política ótima de longo prazo por meio da experiência acumulada.
  \item Representação de um conjunto de tamanho variável de nós candidatos ao descarregamento em um vetor de estado compacto de seis dimensões na qual a política de controle treinada em um cenário representativo pode ser testada em ambientes inéditos sem a necessidade de retreinamento.
  \item Redução do espaço de ações a uma decisão binária (processamento local vs.\ melhor candidato de descarregamento), possibilitando a inferência em $O(1)$ passagens diretas na rede neural por tarefa, tornando o modelo adequado para a implantação embarcada em tempo real.
  \item Desenvolvimento de uma função de recompensa sensível ao prazo que fornece um sinal de gradiente contínuo em ambos os regimes de sucesso e falha, evitando a esparsidade de recompensa e ensinando o agente a minimizar ativamente o atraso em vez de meramente evitá-lo.
\end{enumerate}

O restante deste artigo está organizado da seguinte forma. A Seção~\ref{sec:related} revisa os trabalhos relacionados presentes na literatura. A Seção~\ref{sec:model} formaliza o modelo do sistema e define o objetivo de otimização. A Seção~\ref{sec:arch} detalha a arquitetura da política híbrida \acrotip{SAC}-\acrotip{TOPSIS} proposta. A Seção~\ref{sec:eval} apresenta a avaliação experimental e discute os resultados obtidos. A Seção~\ref{sec:conclusion} conclui o artigo, apresentando as considerações finais e direções para trabalhos futuros. Além disso, as definições de todas as siglas utilizadas ao longo do texto encontram-se sumarizadas na Tabela~\ref{tab:acronyms}.

\begin{table*}[t]
\caption{Lista de Acrônimos -- Resumo das principais siglas utilizadas ao longo deste artigo.}
\label{tab:acronyms}
{\scriptsize
\begin{tabularx}{\textwidth}{l X l X}
\toprule
\textbf{Sigla} & \textbf{Definição} & \textbf{Sigla} & \textbf{Definição} \\
\midrule
\acrotip{3GPP}   & 3rd Generation Partnership Project & \acrotip{NIS}    & Solução Ideal Negativa \\
\acrotip{5G}     & 5th Generation                      & \acrotip{NRV2X}  & New Radio Vehicle-to-Everything \\
\acrotip{CV2X}   & Cellular Vehicle-to-Everything      & \acrotip{OBUI}   & Unidades de Bordo Interconectadas \\
\acrotip{D3QN}   & Dueling Double Deep Q-Network       & \acrotip{PIS}    & Solução Ideal Positiva \\
\acrotip{DRL}    & Aprendizado por Reforço Profundo    & \acrotip{QoS}    & Qualidade de Serviço \\
\acrotip{DSRC}   & Dedicated Short-Range Comm.         & \acrotip{RSU}    & Roadside Unit \\
\acrotip{IEEE}   & Inst. of Elect. and Electr. Eng.    & \acrotip{SAC}    & Soft Actor-Critic \\
\acrotip{IoT}    & Internet of Things                  & \acrotip{TOPSIS} & Tech. Order Pref. Sim. Ideal Sol. \\
\acrotip{IoV}    & Internet de Veículos                & \acrotip{V2I}    & Vehicle-to-Infrastructure \\
\acrotip{ITS}    & Sistemas Transporte Inteligentes    & \acrotip{V2N}    & Vehicle-to-Network \\
\acrotip{JCT}    & Tempo de Conclusão de Tarefa        & \acrotip{V2P}    & Vehicle-to-Pedestrian \\
\acrotip{MCDM}   & Tomada de Decisão Multicritério     & \acrotip{V2V}    & Vehicle-to-Vehicle \\
\acrotip{MEC}    & Computação de Borda de Múltiplos Acessos           & \acrotip{VEC}    & Computação de Borda Veicular \\
\bottomrule
\end{tabularx}}
\end{table*}

